\documentclass[11pt,a4paper]{article}
\usepackage[margin=23mm]{geometry}
\usepackage{amsmath,amssymb,booktabs,microtype}
\usepackage[hidelinks]{hyperref}
\newcommand{\dd}{\mathrm d}
\newcommand{\one}{\mathbf 1}
\title{Route G-R5: Acquired Clock Data and\protect\\an External-Reference Observation Model}
\author{An empirical eligibility boundary and two measurable contrasts}
\date{24 September 2026}
\begin{document}
\maketitle
\begin{abstract}
Public E2/E3 comparison data were acquired and audited, including
2,324,201 records from the 2025 ROCIT release. They do not include the
paired bath-state, independent radiometry and correction ledger needed
for the requested two-temperature test. That empirical task remains
open. Separately, we specify a second Yb$^+$ E3 clock outside the changed
environment and derive both differential and common target-reference
observables. A thermal-response-cancelling composite is useful, but
reference drift and common instrumental shifts remain exactly degenerate
with an unexplained common signal unless independently constrained.
No universal time law is fitted and no new experimental effect is claimed.
\end{abstract}

\section{What the acquired data can and cannot establish}
Two primary releases were obtained with checksums and source attribution:
\begin{itemize}
\item PTB's 2021 release~\cite{ptb} contains 11 published E3/E2 ratio
points, their times and uncertainty components. Its ratio column is
an \emph{absolute ratio offset}, not a fractional offset. The associated
campaign used two clock systems. No paired thermal calibration is supplied.
\item ROCIT's September 2025 release~\cite{rocit} contains 34 selected
same-package E2/E3 comparator files from the 2022 campaign: 2,324,201
rows, all marked valid, no nonfinite values or duplicate timestamps.
Its columns are time, comparator output, validity and two systematic
uncertainties. It supplies neither two-temperature assignments nor
independent ion-site radiometry and a per-run BBR correction ledger.
\end{itemize}
Original release bytes are not the same as uncorrected atom/servo data.
The complete small PTB archive and one exact ROCIT day plus metadata
are preserved; every selected ROCIT member is hashed in the full audit.
No thermal slope is extracted from these unsuitable records. This is
a limitation of the audited releases, not a claim that qualifying data
exist nowhere. The later public ROCIT deposit also supersedes the
paper's earlier generic ``on request'' access statement~\cite{network}.

The ROCIT format~\cite{format} defines
$\Delta_{A\to B}=(\widehat\nu_B-\rho^0_{BA}\widehat\nu_A)/s_B$.
For this pair $A$ is E3 and $B$ E2, opposite to our preferred E3/E2
ratio. Exact conversion requires its reference convention, not merely
negating a column. Valid records need not be independent samples.

\newpage
\section{Specify the outside reference, including its response}
Let $S$ be one target Yb$^+$ ion, interleaved on E2 and E3 inside a
separately controlled radiation enclosure. Let $R$ be a second,
independent Yb$^+$ E3 clock in a separate monitored enclosure outside
the changed bath. Compare both target outputs to $R$ using a frequency
comb and phase-stabilized optical paths. This is a specified arrangement,
not hardware constructed here. Existing separate-clock and external-clock
comparisons support the architecture~\cite{network,steinel}, not our
proposed thermal experiment's performance.

Retain G-R4's ordinary electric-dipole response at \emph{both} sites:
\begin{equation}\label{eq:b}
 b_i^X=-\frac{1}{2h\epsilon_0\nu_i^0}\,\mathcal P
 \int_0^\infty\Delta\alpha_i(\omega)u_\omega^X\dd\omega,
 \qquad B_i^X=\log(1+b_i^X),\quad X=S,R.
\end{equation}
For ideal thermal radiation $u^X=a_{\rm rad}T_X^4$; at leading static
order $b_i^X=\beta_i u^X$. A real multi-surface bath generally needs its
spectrum and geometry, not one unexamined wall temperature. Motional
sideband thermometry does not measure radiation temperature.

Allow a candidate common local factor $F_X>0$ without choosing a law
for it. Write the measured logarithmic ratio as
\begin{align}
 q_i&=\log\left(\frac{\nu_i^S/\nu_3^R}{\nu_i^0/\nu_3^0}\right),\\
 q_i&=\chi_S-\chi_R+B_i^S-B_3^R
       +g_{SR}+\ell_i+A_i^S-A_3^R+n_i,\qquad \chi_X=\log F_X.
       \label{eq:observation}
\end{align}
$A$ denotes the remaining ordinary \emph{log-frequency} corrections,
$\ell_i$ the path/comb contribution and $n_i$ measurement noise.
Using logs makes multiplicative factors additive; at first order $A$
is the usual fractional nuisance shift. No reference term is set to
zero simply because the reference is outside the enclosure. In a
weak-field, slow-motion convention one ordinary contribution is
\begin{equation}
 g_{SR}=\frac{\Phi_S-\Phi_R}{c^2}
       -\frac{v_S^2-v_R^2}{2c^2},
\end{equation}
with additional transfer effects assigned to $\ell_i$ to avoid double
counting. Fixed heights remove a static offset, not an unknown changing
one. Conventional spacetime is assumed; no identification with Chen's
$E_d$ or completion of the finite G-R1--G-R3 constraints is implied.

For hot-minus-cold contrast $\Delta$, independently correcting the
ordinary terms yields
\begin{equation}\label{eq:residual}
 z_i=\Delta q_i-\Delta B_i^S+\Delta B_3^R
       -\Delta g_{SR}-\Delta\ell_i-\Delta A_i^S+\Delta A_3^R
     =\Delta(\chi_S-\chi_R)+\Delta n_i .
\end{equation}
Real measurements average with sensitivity functions $W_i(t)$.
Reference noise cancels in $q_3-q_2$ only for matched effective windows:
otherwise it leaves $-\int(W_3-W_2)q_R(t)\dd t$. Interleaving alone
does not guarantee cancellation; dead times and servo response matter.

\newpage
\section{Two contrasts and an exact identifiability obstruction}
Suppress already evaluated ordinary terms for clarity. To first order,
the two external comparisons have the structure
\begin{equation}\label{eq:design}
 \begin{pmatrix}y_2\\y_3\end{pmatrix}
 =\begin{pmatrix}1&\beta_2\\1&\beta_3\end{pmatrix}
   \begin{pmatrix}\chi\\\Delta u_S\end{pmatrix}
   -\beta_3\Delta u_R\one+c\one+\boldsymbol\epsilon,
 \qquad \chi=\Delta(\chi_S-\chi_R).
\end{equation}
Here $c$ represents an unevaluated common path, gravity or reference
drift term. The differential observable is
\begin{equation}
 D=y_3-y_2=(\beta_3-\beta_2)\Delta u_S+\epsilon_3-\epsilon_2.
\end{equation}
Both $\chi$ and a shared reference term cancel. The composite
\begin{align}
 C&=w_2y_2+w_3y_3,\\
 w_2&=\frac{\beta_3}{\beta_3-\beta_2},\qquad
 w_3=-\frac{\beta_2}{\beta_3-\beta_2},\qquad
 w_2+w_3=1,\quad w_2\beta_2+w_3\beta_3=0
\end{align}
instead gives $C=\chi-\beta_3\Delta u_R+c+w^T\boldsymbol\epsilon$.
With G-R4's independently pinned inputs,
\begin{equation}
 (w_2,w_3)\simeq(-0.16004,1.16004).
\end{equation}
This cancels the leading target $T^4$ response, not all dynamic or
other shifts. Negative weighting can increase noise. The design matrix
has determinant $\beta_3-\beta_2\ne0$: with reference effects controlled,
two external channels can distinguish a common change from the selected
thermal response. But solving two channels for two unknowns leaves
\emph{no residual test}. Independent radiometry is needed to test the
thermal law, rather than infer the bath from that law and call it verified.

\textbf{Obstruction.} If $c$ is unconstrained, the transformation
$(\chi,c)\mapsto(\chi+\lambda,c-\lambda)$ leaves all observations
unchanged. Thus even this external comparison measures an unexplained
common residual, not uniquely a time effect. A factor common to both
environments cancels as before. A third ratio constructed from the same
two channels supplies only a closure identity, not another independent
equation; an independently measured link can diagnose instrumental closure.

For independently calibrated corrections in Eq.~\eqref{eq:residual},
the generalized least-squares common estimate and its variance are
\begin{equation}
 \widehat\chi=\frac{\one^TV^{-1}z}{\one^TV^{-1}\one},\qquad
 \operatorname{Var}\widehat\chi=(\one^TV^{-1}\one)^{-1}.
\end{equation}
This is an estimator specification, not a fit performed here. Shared
reference error contributes $\sigma_R^2\one\one^T$ to $V$; it cancels
from the differential variance but remains a floor in the common estimate.
Correction covariance and sensitivity to the same calibration inputs
must also be propagated, rather than counted as independent noise.

\newpage
\section{Numerical control requirements and the empirical handoff}
For the illustrative target change $300\to310\,\mathrm K$, G-R4's
static coefficients predict
\begin{equation}
 (\Delta b_2^S,\Delta b_3^S)\simeq
 (-7.337\,10^{-17},-1.012\,10^{-17}),\qquad
 D\simeq6.325\,10^{-17}.
\end{equation}
These are calculations, not observations. With fixed nominal composite
weights, the uncertainty from independently calibrated polarizabilities is
\begin{equation}
 \delta C=\Delta u_S(w_2\delta\beta_2+w_3\delta\beta_3),\quad
 \sigma_C\in[|a_2-a_3|,a_2+a_3],\quad
 a_i=|w_i|\sigma_{\beta_i}|\Delta u_S|.
\end{equation}
Unknown correlation gives $\sigma_C$ between $2.17\,10^{-18}$ and
$2.59\,10^{-18}$ from these inputs alone. This is not a confidence
interval or an achieved total uncertainty. The retained historical
calibrations do not support a sub-$10^{-18}$ common-mode claim.
For scale, a reference bath change $300\to300.1\,\mathrm K$ contributes
about $-9.64\,10^{-20}$ to its E3 frequency, while a $1\,\mathrm{cm}$
height change gives $g\Delta h/c^2\simeq1.09\,10^{-18}$. These are
conditional sensitivities, not asserted apparatus fluctuations.

Use settled low--high--high--low blocks with matched effective centers
$t=(-3,-1,1,3)\tau$ and contrast weights $(-1,1,1,-1)/2$.
Then $\sum w=\sum wt=0$ and the hot-minus-cold response has unit gain.
However $\sum wt^2=-8\tau^2$: nonlinear drift survives, as does
heater-correlated bias. Reversed ordering, monitored reference
cross-talk and a dummy-heater/electrical-pickup control are useful
discriminators, not proofs of cancellation. No particular dwell time
or hardware precision is asserted without the thermal/servo response.

The next empirical input must supply uncorrected comparator/servo
records, bath-state assignments, independent radiation calibration at
both sites, every applied correction and its covariance, synchronized
validity/averaging windows, and reference/path/gravity controls. The
machine-readable contract is \texttt{EXTERNAL\_REFERENCE\_CARD.json}.
No researcher has been contacted or new experiment performed. Obtaining
those records remains open; the acquisition audit and reference-model
derivation are bounded-done. The 46 new checks verify algebra and source
eligibility, not evidence for a common time effect.

\begin{thebibliography}{9}\small
\bibitem{ptb} R. Lange et al., PTB comparison data (2021),
\href{https://doi.org/10.7795/720.20211028}{doi:10.7795/720.20211028};
associated paper: \href{https://doi.org/10.1103/PhysRevLett.126.011102}{Phys. Rev. Lett. 126, 011102}.
\bibitem{rocit} T. Lindvall et al., ROCIT campaign dataset (2025),
\href{https://doi.org/10.5281/zenodo.17107693}{doi:10.5281/zenodo.17107693}.
\bibitem{network} T. Lindvall et al., \emph{Coordinated international
comparisons between optical clocks connected via fiber and satellite
links}, \href{https://doi.org/10.1364/OPTICA.561754}{Optica 12, 843 (2025)}.
\bibitem{format} INRIM, \emph{Optical link data format},
\href{https://github.com/INRIM/optical-link-data-format}{source specification}, accessed 24 September 2026.
\bibitem{steinel} M. Steinel et al., \emph{Evaluation of a $^{88}$Sr$^+$
Optical Clock with a Direct Measurement of the Blackbody Radiation
Shift and Determination of the Clock Frequency},
\href{https://doi.org/10.1103/PhysRevLett.131.083002}{Phys. Rev. Lett. 131, 083002 (2023)}.
\end{thebibliography}
\end{document}
